% Źródła: materialy/2026-10-12/notatki.pdf i zdjęcia IMG_6917--IMG_6922.
\chapter{Prawdopodobieństwo warunkowe}
\label{chap:warunkowe}

\section{Definicja i podstawowe własności}
\label{sec:warunkowe}

Niech $(\Omega,\F,\PP)$ będzie przestrzenią probabilistyczną.

\begin{definicja}
\label{def:warunkowe}
Dla zdarzeń $A,B\in\F$ takich, że $\PP(B)>0$,
\emph{prawdopodobieństwo warunkowe zdarzenia $A$ pod warunkiem $B$}
definiujemy wzorem
\begin{equation}
  \PP(A\mid B)=\frac{\PP(A\cap B)}{\PP(B)}.
  \label{eq:warunkowe}
\end{equation}
\end{definicja}

\begin{uwaga}
Przy ustalonym $B\in\F$ o dodatnim prawdopodobieństwie wzór
\[
  \PP_B(A)=\PP(A\mid B),\qquad A\in\F,
\]
definiuje nowe prawdopodobieństwo na $(\Omega,\F)$.
\end{uwaga}

\begin{fakt}
\label{fakt:wlasnosci-warunkowego}
Niech $A,B\in\F$ oraz $\PP(B)>0$. Wówczas:
\begin{enumerate}
  \item jeśli $A\subseteq B$, to
    \[
      \PP(A\mid B)=\frac{\PP(A)}{\PP(B)};
    \]
  \item jeśli $B\subseteq A$, to $\PP(A\mid B)=1$;
  \item jeśli $A\cap B=\varnothing$, to $\PP(A\mid B)=0$;
    ogólniej, wystarczy warunek $\PP(A\cap B)=0$;
  \item jeśli $C\in\F$ i $\PP(B\cap C)>0$, to
    \begin{equation}
      \PP(A\mid B\cap C)
        =\frac{\PP(A\cap B\mid C)}{\PP(B\mid C)}.
      \label{eq:warunkowanie-dwukrotne}
    \end{equation}
\end{enumerate}
\end{fakt}

\begin{przyklad}[Rodzina z dwojgiem dzieci]
\label{przykl:dzieci}
Wybieramy rodzinę z dwojgiem dzieci. Jakie jest prawdopodobieństwo,
że ma ona dwóch chłopców, jeśli:
\begin{enumerate}[label=\alph*)]
  \item starsze dziecko jest chłopcem;
  \item w rodzinie jest co najmniej jeden chłopiec?
\end{enumerate}
Zakładamy, że płcie dzieci są niezależne, a każda z dwóch płci ma
prawdopodobieństwo $1/2$.
\end{przyklad}

\noindent\emph{Rozwiązanie.}
Kodujemy chłopca przez $0$, a dziewczynkę przez $1$. Pierwsza współrzędna
odpowiada starszemu dziecku. Cztery stany elementarne
\[
  \Omega=\{(0,0),(0,1),(1,0),(1,1)\}
\]
są jednakowo prawdopodobne. Przyjmujemy
\begin{align*}
  A&=\{(0,0)\},\\
  B_1&=\{(0,0),(0,1)\},\\
  B_2&=\{(0,0),(0,1),(1,0)\}.
\end{align*}
Ponieważ $A\subseteq B_1$ i $A\subseteq B_2$, otrzymujemy
\[
  \PP(A\mid B_1)=\frac{1/4}{1/2}=\frac12,
  \qquad
  \PP(A\mid B_2)=\frac{1/4}{3/4}=\frac13.
\]

\section{Prawdopodobieństwo całkowite}
\label{sec:calkowite}

\begin{definicja}
\label{def:rozbicie}
\emph{Rozbiciem przestrzeni $\Omega$} nazywamy rodzinę zbiorów
$(B_n)_{n\in I}$, gdzie $I\subseteq\mathbb{N}$ jest zbiorem indeksów,
taką, że:
\begin{enumerate}
  \item $B_n\cap B_m=\varnothing$ dla $n,m\in I$, $n\ne m$;
  \item $\displaystyle\bigcup_{n\in I}B_n=\Omega$.
\end{enumerate}
Rozbicie może być skończone lub przeliczalne.
\end{definicja}

\begin{twierdzenie}[O prawdopodobieństwie całkowitym]
\label{tw:calkowite}
Niech $(B_n)_{n\in I}$ będzie rozbiciem $\Omega$ takim, że
$B_n\in\F$ i $\PP(B_n)>0$ dla każdego $n\in I$.
Wówczas dla dowolnego $A\in\F$ zachodzi
\begin{equation}
  \PP(A)=\sum_{n\in I}\PP(A\mid B_n)\PP(B_n).
  \label{eq:calkowite}
\end{equation}
\end{twierdzenie}

\begin{proof}
Zdarzenia $A\cap B_n$ są parami rozłączne. Zatem
\begin{align*}
  \PP(A)
    &=\PP(A\cap\Omega)
      =\PP\!\left(A\cap\bigcup_{n\in I}B_n\right)\\
    &=\PP\!\left(\bigcup_{n\in I}(A\cap B_n)\right)
      =\sum_{n\in I}\PP(A\cap B_n)\\
    &=\sum_{n\in I}\PP(A\mid B_n)\PP(B_n).
\end{align*}
\end{proof}

\begin{przyklad}[Zniszczenie celu]
\label{przykl:strzaly}
Prawdopodobieństwo trafienia w cel wynosi $1/2$, a wyniki kolejnych
strzałów są niezależne. Prawdopodobieństwo zniszczenia celu przy
dokładnie $k$ celnych strzałach wynosi $1-(1/3)^k$.
Jakie jest prawdopodobieństwo zniszczenia celu po oddaniu dziesięciu strzałów?
\end{przyklad}

\noindent\emph{Rozwiązanie.}
Niech $A$ oznacza zniszczenie celu, a $B_i$ zdarzenie, że dokładnie
$i$ spośród dziesięciu strzałów jest celnych, gdzie $i=0,1,\ldots,10$.
Wówczas
\[
  \PP(A\mid B_i)=1-\left(\frac13\right)^i,
  \qquad
  \PP(B_i)=\binom{10}{i}\left(\frac12\right)^{10}.
\]
Korzystając ze wzoru \eqref{eq:calkowite} i ze wzoru dwumianowego
\[
  (a+b)^n=\sum_{i=0}^{n}\binom{n}{i}a^ib^{n-i},
\]
otrzymujemy
\begin{align*}
  \PP(A)
    &=\sum_{i=0}^{10}\left(1-\left(\frac13\right)^i\right)
        \binom{10}{i}\left(\frac12\right)^{10}\\
    &=\sum_{i=0}^{10}\binom{10}{i}\left(\frac12\right)^{10}
      -\sum_{i=0}^{10}\binom{10}{i}
        \left(\frac16\right)^i\left(\frac12\right)^{10-i}\\
    &=1-\left(\frac16+\frac12\right)^{10}
      =\boxed{1-\left(\frac23\right)^{10}}.
\end{align*}

\begin{przyklad}[Dwie urny]
\label{przykl:urny}
Mamy dwie urny, każda zawiera sto kul. Jedna kula jest biała:
trafia do pierwszej urny z prawdopodobieństwem $2/3$, a do drugiej
z prawdopodobieństwem $1/3$. Pozostałe kule są czarne.
Wykonujemy osiemdziesiąt losowań ze zwracaniem.
Jak rozdzielić losowania między urny, aby zmaksymalizować
prawdopodobieństwo wylosowania białej kuli co najmniej raz?
\end{przyklad}

\noindent\emph{Rozwiązanie.}
Niech $B_i$ oznacza zdarzenie, że biała kula znajduje się w $i$-tej urnie,
a $A$ zdarzenie, że zostanie wylosowana co najmniej raz.
Wykonujemy $k$ losowań z pierwszej urny i $m=80-k$ z drugiej.
Oznaczmy przez $q=99/100$ prawdopodobieństwo wylosowania czarnej kuli
z urny zawierającej białą kulę. Wówczas
\[
  \PP(A\mid B_1)=1-q^k,
  \qquad
  \PP(A\mid B_2)=1-q^{80-k}.
\]
Ze wzoru \eqref{eq:calkowite} otrzymujemy
\begin{equation}
  \PP(A)=f(k)=\frac23(1-q^k)+\frac13(1-q^{80-k}),
  \qquad k\in\{0,1,\ldots,80\}.
  \label{eq:urny-funkcja}
\end{equation}
Aby znaleźć maksimum, rozważamy najpierw zmienną rzeczywistą $x\in[0,80]$.
Liczymy
\[
  f'(x)=\frac{\ln q}{3}\bigl(q^{80-x}-2q^x\bigr).
\]
Warunek $f'(x)=0$ jest równoważny
\[
  q^{80-2x}=2,
  \qquad
  x=\frac12\left(80-\frac{\ln2}{\ln q}\right)\approx74{,}484.
\]
Ponieważ
\[
  f''(x)=-\frac{(\ln q)^2}{3}\bigl(q^{80-x}+2q^x\bigr)<0,
\]
jest to jedyny punkt maksimum. Dla całkowitej liczby losowań
porównujemy dwóch kandydatów:
\[
  f(74)\approx0{,}3692799222,
  \qquad f(75)\approx0{,}3692788890.
\]
Najlepiej wykonać \emph{74 losowania z pierwszej urny i 6 z drugiej}.
Oba powyższe prawdopodobieństwa po zaokrągleniu do czterech miejsc
po przecinku wynoszą $0{,}3693$.

\section{Wzór Bayesa}
\label{sec:bayes}

\begin{twierdzenie}[Wzór Bayesa]
\label{tw:bayes}
Niech $(B_n)_{n\in I}$ będzie rozbiciem $\Omega$ takim, że
$B_n\in\F$ i $\PP(B_n)>0$ dla każdego $n\in I$.
Jeżeli $A\in\F$ i $\PP(A)>0$, to dla każdego $n\in I$
\begin{equation}
  \PP(B_n\mid A)
    =\frac{\PP(A\mid B_n)\PP(B_n)}{\PP(A)}
    =\frac{\PP(A\mid B_n)\PP(B_n)}
      {\displaystyle\sum_{k\in I}\PP(A\mid B_k)\PP(B_k)}.
  \label{eq:bayes}
\end{equation}
\end{twierdzenie}

\begin{przyklad}[ZONK]
\label{przykl:zonk}
Za jedną z trzech bramek znajduje się nagroda, a za pozostałymi dwoma
zonki. Każda bramka zawiera nagrodę z prawdopodobieństwem $1/3$.
Wybieramy pierwszą bramkę. Filip Chajzer zna położenie nagrody
i otwiera jedną z pozostałych bramek, zawsze pokazując zonka.
Jeżeli może otworzyć obie, wybiera każdą z prawdopodobieństwem $1/2$.
Chajzer otworzył trzecią bramkę. Jakie jest prawdopodobieństwo,
że nagroda znajduje się w drugiej bramce? Czy warto zmienić wybór?
\end{przyklad}

\noindent\emph{Rozwiązanie.}
Przyjmujemy oznaczenia
\[
  A=\{\text{Chajzer otwiera trzecią bramkę}\},
  \qquad
  B_i=\{\text{nagroda jest w $i$-tej bramce}\},
\]
gdzie $i=1,2,3$. Zdarzenia $B_1,B_2,B_3$ tworzą rozbicie $\Omega$,
a $\PP(B_i)=1/3$ dla każdego $i$.
Jeśli nagroda znajduje się w pierwszej bramce, prowadzący ma wybór
między drugą a trzecią. Jeśli znajduje się w drugiej, musi otworzyć
trzecią; jeśli w trzeciej, musi otworzyć drugą. Stąd
\[
  \PP(A\mid B_1)=\frac12,
  \qquad \PP(A\mid B_2)=1,
  \qquad \PP(A\mid B_3)=0.
\]
Ze wzoru na prawdopodobieństwo całkowite \eqref{eq:calkowite} mamy
\[
  \PP(A)=\frac12\cdot\frac13+1\cdot\frac13+0\cdot\frac13=\frac12.
\]
Zatem wzór Bayesa \eqref{eq:bayes} daje
\begin{equation}
  \PP(B_2\mid A)
    =\frac{\PP(A\mid B_2)\PP(B_2)}{\PP(A)}
    =\frac{1\cdot(1/3)}{1/2}=\boxed{\frac23}.
  \label{eq:zonk-wynik}
\end{equation}
Natomiast $\PP(B_1\mid A)=1/3$. Warto więc zmienić wybór z pierwszej
bramki na drugą: prawdopodobieństwo wygranej wzrasta z $1/3$ do $2/3$.

\begin{przyklad}[Moneta z dwoma orłami]
\label{przykl:moneta}
Mamy sto monet: dziewięćdziesiąt dziewięć uczciwych i jedną
z orłem na obu stronach. Wybieramy monetę losowo i rzucamy nią dziesięć razy.
Za każdym razem wypada orzeł. Jakie jest prawdopodobieństwo,
że wybraliśmy monetę z dwoma orłami?
\end{przyklad}

\noindent\emph{Rozwiązanie.}
Niech $A$ oznacza dziesięć orłów, $B_1$ wybór uczciwej monety,
a $B_2$ wybór monety z dwoma orłami. Wówczas
\begin{align*}
  \PP(B_1)&=\frac{99}{100},&
  \PP(B_2)&=\frac1{100},\\
  \PP(A\mid B_1)&=\left(\frac12\right)^{10},&
  \PP(A\mid B_2)&=1.
\end{align*}
Prawdopodobieństwo dziesięciu orłów wynosi
\[
  \PP(A)=\left(\frac12\right)^{10}\frac{99}{100}+\frac1{100}.
\]
Ze wzoru Bayesa \eqref{eq:bayes} otrzymujemy
\begin{equation}
  \PP(B_2\mid A)
    =\frac{1/100}{(1/2)^{10}\cdot99/100+1/100}
    =\frac{1024}{1123}\approx\boxed{0{,}9118}.
  \label{eq:moneta-wynik}
\end{equation}
